\documentclass[11pt]{article}

\usepackage{nmrdoc}
\usepackage{siunitx}
\usepackage{bm}

\sisetup{per-mode=symbol}
\DeclareSIUnit{\angstrom}{\text{\AA}}

\begin{document}
\NmrTitle{The RingSusceptibility Calculator}

\section{What the calculator computes}

The RingSusceptibility calculator describes the through-space magnetic effect
of an aromatic ring treated as a small anisotropic magnetic object at the ring
centre. In a Pople-style ring-current susceptibility model, the applied NMR
field polarizes the delocalized \(\pi\) electrons of the ring. Their response
produces a secondary magnetic field at a nearby nucleus, and that secondary
field changes the magnetic shielding felt by the nucleus.

The implementation keeps the geometric part of this idea. For each target atom
and each accepted aromatic ring, it builds a \(3\times 3\) tensor from the
ring-centre position, the ring normal, and the vector from the ring centre to
the atom. The tensor has units \(\si{\angstrom^{-3}}\). A later calibration can
learn how this descriptor correlates with a DFT shielding contribution. Before
that calibration, the stored number is not a chemical shift and not a shielding
in ppm. The observable to keep in view is the nuclear shielding tensor; in
solution NMR its isotropic part is what enters the reported chemical shift after
reference subtraction and the usual sign convention.

\section{Where shielding enters}

For a nucleus in an applied magnetic field \(\bm B_0\), the local field is
written
\begin{equation}
  \bm B_{\mathrm{nuc}} = (\bm I-\bm\sigma)\bm B_0 .
  \label{eq:shielding}
\end{equation}
Here \(\bm I\) is the identity matrix and \(\bm\sigma\) is the nuclear shielding
tensor. In physical NMR \(\bm\sigma\) is dimensionless: it is the linear
coefficient that tells how the electron cloud changes the field at the nucleus
when \(\bm B_0\) is applied. It is a matrix rather than a scalar because the
molecule has orientation. A field applied along one molecular direction can
induce a secondary field with a different size, or even a different direction,
from a field applied along another direction.

The quantum-mechanical origin is the electronic response to the magnetic field.
The applied field perturbs the electronic state; the perturbed electronic
current and magnetization create a secondary field at the nucleus; the linear
part of that mapping is \(\bm\sigma\). The RingSusceptibility calculator does
not solve that response. It uses the classical far-field shape expected from an
anisotropic ring susceptibility: the ring is represented by its centre and by a
unit normal \(\hat{\bm n}\), the axis perpendicular to the ring plane.

This is the same rank-2 setting as Eq.~\eqref{eq:shielding}. The tensor entry
\(X_{ab}\) says how a field component associated with direction \(b\) contributes
to the induced field component along direction \(a\). The code's tensor is not
yet the physical shielding tensor \(\bm\sigma\), but it has the same
transformation law under rotations. If the molecule is rotated, the \(3\times3\)
array must rotate with it; a scalar distance-to-ring feature would lose that
orientation information.

For an aromatic ring, the direction \(\hat{\bm n}\) is the natural anisotropy
axis. Above and below the ring plane the geometric scalar contribution has one
sign; in the ring plane it has the opposite sign. The angular function behind
that statement is the dipolar factor \(3\cos^2\theta-1\), where \(\theta\) is
the angle between the ring normal and the ring-centre-to-atom direction. This is
the same second-Legendre-polynomial pattern as a magnetic dipole field: two
lobes along the axis, a belt of the opposite sign in the equatorial plane, and a
zero on the cone with \(\cos^2\theta=1/3\), the magic angle near
\(54.7^\circ\). In this calculator that picture is not decoration; it is the
scalar \(T_0\) part of the tensor derived below.

\section{The tensor split used by the code}

The project stores each \(3\times3\) tensor through the split implemented by
\texttt{SphericalTensor::Decompose}. For a real tensor \(\bm X\), the scalar
piece is
\begin{equation}
  T_0 = \frac{1}{3}\operatorname{tr}\bm X .
  \label{eq:t0}
\end{equation}
This is the isotropic component. It multiplies the identity matrix, survives
molecular tumbling, and is the part that contributes to an isotropic shielding
after the geometric descriptor has been calibrated to a physical scale.

The antisymmetric part is stored as the three-component vector
\begin{equation}
  \bm T_1 =
  \frac{1}{2}
  \begin{pmatrix}
    X_{yz}-X_{zy}\\
    X_{zx}-X_{xz}\\
    X_{xy}-X_{yx}
  \end{pmatrix}.
  \label{eq:t1}
\end{equation}
This vector is the dual representation of the antisymmetric matrix part. It is
usually not observed directly in standard isotropic solution NMR, but the ring
susceptibility tensor can be asymmetric, so the code keeps it.

The rest is the symmetric traceless matrix
\begin{equation}
  \bm S =
  \frac{\bm X+\bm X^{\mathsf T}}{2} - T_0\bm I,
  \qquad \operatorname{tr}\bm S=0 .
  \label{eq:symtr}
\end{equation}
Symmetric means \(S_{ab}=S_{ba}\). Traceless means the three diagonal entries
sum to zero. Those two conditions leave five independent numbers, the
orientation-dependent anisotropy. The code packs them as
\begin{equation}
  \bm T_2 =
  \begin{pmatrix}
    \sqrt{2}\,S_{xy}\\
    \sqrt{2}\,S_{yz}\\
    \sqrt{3/2}\,S_{zz}\\
    \sqrt{2}\,S_{xz}\\
    (S_{xx}-S_{yy})/\sqrt{2}
  \end{pmatrix}.
  \label{eq:t2}
\end{equation}
The factors in Eq.~\eqref{eq:t2} are the implementation's real spherical basis.
They make the five-vector length match the Frobenius length of
\(\bm S\): \(\sum_m T_{2m}^2=\sum_{ab}S_{ab}^2\). The stored nine-number order is
\[
  [T_0,T_{1x},T_{1y},T_{1z},
    T_{2,-2},T_{2,-1},T_{2,0},T_{2,+1},T_{2,+2}] .
\]
RingSusceptibility populates all three parts. Its \(T_0\) is not discarded, its
\(T_1\) can be nonzero, and its \(T_2\) carries the angular anisotropy that a
scalar ring-current descriptor would hide.

\section{The method implemented}

The geometry result first assigns each aromatic ring a centre, a normal, and an
average radius. For a ring with vertex positions \(\bm q_k\), the centre is
\begin{equation}
  \bm c = \frac{1}{N}\sum_{k=1}^{N}\bm q_k ,
  \label{eq:center}
\end{equation}
with positions in \(\si{\angstrom}\). The normal \(\hat{\bm n}\) is the unit
vector perpendicular to the best-fit plane of the centred vertex coordinates,
found by an SVD. The sign is chosen from the stored cyclic order of the ring
vertices, although the tensor below is unchanged if \(\hat{\bm n}\) is replaced
by \(-\hat{\bm n}\). The radius \(R\) is the average distance of the vertices
from \(\bm c\).

For a target atom at position \(\bm r_i\), define
\begin{equation}
  \bm x = \bm r_i-\bm c,\qquad
  r = |\bm x|,\qquad
  \hat{\bm d} = \frac{\bm x}{r},\qquad
  u = \hat{\bm d}\cdot\hat{\bm n}.
  \label{eq:ring-geometry}
\end{equation}
The distance \(r\) is in \(\si{\angstrom}\). The vectors \(\hat{\bm d}\) and
\(\hat{\bm n}\) are unitless unit vectors. The scalar \(u\) is the cosine of the
angle between the atom direction and the ring normal.

The per-ring scalar recorded by the code is
\begin{equation}
  f = \frac{3u^2-1}{r^3},
  \label{eq:f}
\end{equation}
with units \(\si{\angstrom^{-3}}\). Equation~\eqref{eq:f} is positive on the
ring normal, negative in the ring plane, and zero at the magic angle. The helper
also forms the symmetric traceless dipolar tensor
\begin{equation}
  K_{ab} =
  \frac{3\hat d_a\hat d_b-\delta_{ab}}{r^3},
  \label{eq:dipolar}
\end{equation}
where \(a,b\in\{x,y,z\}\) and \(\delta_{ab}\) is one when \(a=b\) and zero
otherwise. \(K_{ab}\) is retained inside the helper result, but the per-ring
record and the per-atom feature output use the full tensor
\begin{equation}
  M_{ab} =
  \frac{
    9u\,\hat d_a\hat n_b
    -3\hat n_a\hat n_b
    -\left(3\hat d_a\hat d_b-\delta_{ab}\right)}
  {r^3}.
  \label{eq:full}
\end{equation}
Every numerator term in Eq.~\eqref{eq:full} is unitless, so \(M_{ab}\) has
units \(\si{\angstrom^{-3}}\). The first term, \(\hat{\bm d}\otimes\hat{\bm n}\),
is generally asymmetric and can contribute to \(T_1\). The second term projects
onto the ring-normal axis. The last term is the negative of the dipolar tensor
in Eq.~\eqref{eq:dipolar}.

The trace of Eq.~\eqref{eq:full} gives the code's scalar identity:
\begin{equation}
  \frac{1}{3}\operatorname{tr}\bm M
  = \frac{3u^2-1}{r^3}
  = f .
  \label{eq:t0-is-f}
\end{equation}
The ring scalar is therefore the \(T_0\) component of the full tensor, not a
separate scalar model replacing the tensor.

A textbook susceptibility form would next attach a ring scale factor, for
example
\begin{equation}
  \sigma^{(R)}_{ab} =
  \frac{\Delta\chi_R}{3} M_{ab},
  \label{eq:textbook}
\end{equation}
where \(\Delta\chi_R\) is the anisotropic magnetic susceptibility assigned to
ring \(R\), carrying whatever physical units are needed to turn the
\(\si{\angstrom^{-3}}\) geometry into a dimensionless shielding coefficient.
Equation~\eqref{eq:textbook} is not evaluated by this calculator. The source
code does not read a \(\Delta\chi_R\) value, does not call the ring-current
intensity accessors, and does not use the Johnson--Bovey lobe offsets. The TOML
file contains ring-current intensities in \(\si{\nano\ampere\per\tesla}\)
(PHE \(-12.0\), TYR \(-11.28\), TRP benzene \(-12.48\), TRP pyrrole \(-6.72\),
HIS/HID/HIE \(-5.16\), and TRP indole perimeter \(-19.2\)) and lobe offsets in
\(\si{\angstrom}\), but they belong to other ring-current calculations and ring
metadata. In this class, all accepted aromatic ring types use Eq.~\eqref{eq:full}
with unit geometric weight.

The candidate rings are the aromatic rings exposed by the protein topology:
PHE benzene, TYR phenol, TRP benzene, TRP pyrrole, TRP indole perimeter, and the
HIS/HID/HIE imidazole variants. The saturated proline pyrrolidine ring is stored
elsewhere in the topology and is not returned by the aromatic-ring accessor used
here. A spatial index finds ring centres within \(\SI{15.0}{\angstrom}\) of the
target atom. That cutoff comes from the ring-current spatial cutoff parameter.

Before a candidate ring contributes, three filters must accept it. The
singularity guard rejects distances below \(\SI{0.1}{\angstrom}\), where the
\(r^{-3}\) expression would diverge numerically. The near-field filter uses
the ring diameter \(2R\) as the source extent and the dimensionless ratio
\(0.5\), so it accepts when \(r>R\). This keeps the point-source
susceptibility expression outside the finite ring source. The bonded-ring
filter rejects a target atom if it is a ring vertex or is covalently bonded to
a ring vertex, because those atoms are in a through-bond regime rather than the
through-space regime represented by Eq.~\eqref{eq:full}.

For atom \(i\), the accepted ring tensors are accumulated as
\begin{equation}
  \bm M_i = \sum_{R\in\mathcal A_i} \bm M^{(iR)} ,
  \label{eq:sum}
\end{equation}
where \(\mathcal A_i\) is the set of accepted nearby aromatic rings. The total
\(\bm M_i\) is then decomposed by Eqs.~\eqref{eq:t0}--\eqref{eq:t2}. The
per-ring record also stores \(M_{ab}\), its spherical-tensor split, the scalar
\(f\), the centre-to-atom distance, and ring-frame coordinates
\(z=\bm x\cdot\hat{\bm n}\), \(\rho=|\bm x-z\hat{\bm n}|\), and
\(\theta=\operatorname{atan2}(\rho,|z|)\) in radians.

The same formula can be sampled at an arbitrary point in space. That path sums
over aromatic rings, skips points closer than \(\SI{0.1}{\angstrom}\), skips
points with \(r<R\), and skips points farther than \(\SI{15.0}{\angstrom}\).
It does not run the bonded-ring filter, since a grid point has no atom index or
covalent topology.

\section{Inputs and output}

RingSusceptibility consumes coordinates in \(\si{\angstrom}\), typed aromatic
ring topology, ring geometry, and a spatial index over ring centres. It has no
role for MOPAC charges or bond orders, AIMNet2 charges or embeddings, ff14SB
charges, APBS electrostatic fields, or DSSP secondary-structure assignments;
those data feed other calculators or downstream comparisons. ORCA or other DFT
shielding tensors are also outside this calculation and enter later as the
reference for calibration and evaluation.

The calculator's direct feature output is \texttt{ringchi\_shielding.npy}. It
has one row per atom and nine columns packed as
\([T_0,T_{1x},T_{1y},T_{1z},T_{2,-2},T_{2,-1},T_{2,0},T_{2,+1},T_{2,+2}]\).
The values are the decomposition of the summed geometric tensor in
Eq.~\eqref{eq:sum}, with units inherited from \(\bm M_i\):
\(\si{\angstrom^{-3}}\). When the common feature writer emits the shared sparse
ring table, \texttt{ring\_contributions.npy} stores this calculator's per-ring
chi decomposition in zero-based columns 45--53, in the same packed order.
Calibration against a DFT shielding reference is the step that can attach a
physical scale and test correlation with the shielding contribution. The
calculator by itself produces a geometry-derived tensor descriptor, not a
pointwise DFT shielding prediction.

\end{document}
